\documentclass[11pt, a4paper]{article}
\usepackage[utf8]{inputenc}

% Core packages
\usepackage{amsmath}
\usepackage{amssymb}
\usepackage{bm}
\usepackage{geometry}
\usepackage{hyperref}
\usepackage{enumitem}
% \usepackage{cite} % removed: conflicts with biblatex (which handles citation compression natively)
\usepackage{booktabs}
\usepackage{multirow}
\usepackage{makecell}
\usepackage[table]{xcolor}
\usepackage{graphicx}
\usepackage{caption}
\usepackage{tabularx}
\usepackage{titlesec}
\usepackage{setspace}
\usepackage{array}
\usepackage{fancyhdr}

\setlength{\abovecaptionskip}{2pt}
\setlength{\belowcaptionskip}{2pt}
\setlength{\intextsep}{4pt}
\setlength{\textfloatsep}{4pt}
\setlength{\floatsep}{12pt}

\pagestyle{fancy}
\fancyhf{}
\fancyfoot[C]{\thepage}
\renewcommand{\headrulewidth}{0pt}

\setlist[itemize]{itemsep=0pt, parsep=0pt, topsep=0pt, partopsep=0pt}
\setlist[enumerate]{itemsep=0pt, parsep=0pt, topsep=0pt, partopsep=0pt}

\geometry{a4paper, margin=25mm}
\allowdisplaybreaks

\usepackage[backend=biber, style=numeric, sorting=none]{biblatex}
\addbibresource{ref.bib}

\titleformat{\section}{\normalfont\Large\bfseries\setstretch{1.2}}{\thesection}{1em}{}
\titlespacing{\section}{0pt}{4pt}{4pt}
\titleformat{\subsection}{\normalfont\large\bfseries\setstretch{1.1}}{\thesubsection}{1em}{}
\titlespacing{\subsection}{0pt}{4pt}{4pt}
\titleformat{\subsubsection}{\normalfont\normalsize\bfseries\setstretch{1.0}}{\thesubsubsection}{1em}{}
\titlespacing{\subsubsection}{0pt}{4pt}{3pt}

\title{\textbf{QAD-MultiGuard: Multimodal Fraud Detection on Mobile Devices via Quantization-Aware Distillation}}

\author{Author information pending\thanks{Author names, affiliations, and funding information to be added in the final version.}}
\date{}

\begin{document}
\maketitle
\thispagestyle{plain}

% ============================================================================
% Abstract
% ============================================================================
\noindent{\bfseries Abstract:}
Telecommunication fraud poses a critical challenge to mobile edge security, where on-device real-time detection must satisfy three coupled constraints: strict privacy compliance, millisecond-level latency, and resilience to precision degradation under low-bit quantization. We present \textbf{QAD-MultiGuard}, an edge--cloud collaborative, privacy-preserving multimodal fraud detection framework. At its core lies a pure KL-divergence distillation paradigm with a homologous teacher: knowledge from a full-precision (BF16) teacher model is transferred to its quantized students. This paradigm supports a heterogeneous dual-track quantization scheme---server-side native NVFP4 and on-device block-wise GGUF-Q4\_K\_M. To stabilise quantization-sensitive projection layers, we introduce the \textbf{Output-Variance Freeze (OV-Freeze)} regularizer, whose gradient compatibility and joint-loss convergence are formally analysed in Appendix~A. For privacy-preserving multimodal fusion, we construct a 128-dimensional irreversible acoustic embedding ($F_v$) and an L-BFGS-optimised four-modality risk fusion mechanism. On the TeleAntiFraud-28k benchmark, NVFP4 QAD + OV-Freeze attains $F_1 = 0.923 \pm 0.006$, recovering 99.1\% of the BF16 upper bound and outperforming 10 representative baselines (e.g., AWQ, QuaRot, BitDistiller) by 6.5 to 8.5 $F_1$ points. On the in-house adversarial set AdvFraud-3k, the model attains $F_1 = 0.875$, corresponding to a 6.0\% relative drop (absolute drop 0.056) against the BF16 upper bound, and a 5.2\% drop (absolute 0.048) against the deployed model---well within the threat-model constraint $C_2(\le 6\%)$. End-to-end deployment on the Snapdragon~8~Gen~3 platform yields a median alert latency of 268~ms (P99: 342~ms), confirming the framework's viability under real-world anti-fraud workloads.

\vspace{0.5em}
\noindent{\bfseries Keywords:} Telecommunication fraud detection; quantization-aware distillation; privacy preservation; multimodal fusion; speculative decoding; edge--cloud collaboration.

% ============================================================================
% 1. INTRODUCTION
% ============================================================================
\section{Introduction}

The proliferation of generative AI has driven telecommunication fraud from script-based single-modality attacks toward precision-targeted, covert multimodal campaigns. Ma et al.~\cite{1} characterised 14 typical fraud-script patterns and 7 acoustic disguise strategies on the TeleAntiFraud-28k corpus, with each case involving on average 2.7 heterogeneous modality signals including SMS phishing, voice calls, malicious URLs, and social-engineering scripts. Single-modality solutions (e.g., ASR-transcript-based text classifiers) exhibit pronounced detection lag in production deployments. Wang et al.~\cite{2}, reporting on the SAFE-QAQ production system, observed that ASR-transcript-based detection recall drops to 60--75\% under dialect, environmental noise, and deliberate voice deformation interference.

Conventional rule-based or traditional machine learning approaches cannot simultaneously satisfy the real-time, privacy, and on-device resource constraints. First, a typical scam call lasts $\sim$90~s, with critical fraud instructions appearing within the first 60~s. Two distinct latency notions must be distinguished here: (i) the \emph{per-window incremental inference latency}---the time to complete one risk-scoring iteration after each streaming sliding window is triggered---which we require to be $\le 500$~ms so that risk scores refresh at a 1~s cadence (measured P50 $=268$~ms, see \S~3.3.1); and (ii) the \emph{first-alert latency}---the time from call onset to the first valid risk verdict---bounded by the acoustic window length $W$ (here $W=3$~s, with first alert within $\sim$3.3~s). Since critical fraud instructions concentrate in the first 60~s, a 3.3~s first alert plus 1~s continuous refresh provides ample interception opportunities within the coercion window. Second, China's Personal Information Protection Law (PIPL) restricts cross-border transfer of raw audio and other sensitive biometric features, mandating that core processing occur on-device. Third, mobile devices face strict hardware constraints (RAM $\le 4$~GB, available storage $\le 500$~MB). Consequently, deploying a 4-bit quantized large language model on-device emerges as a viable path balancing privacy compliance with real-time responsiveness.

However, standard post-training quantization (PTQ) applied to 0.5B-scale models incurs 6.0\%--12.5\% accuracy degradation on standard benchmarks (per the NVIDIA Nemotron report~\cite{3}); on our telecommunication fraud task (TAF-28k), the $F_1$ gap between PTQ and QAD reaches 6.5--8.5 percentage points (see \S~4.2, Table~\ref{tab3-en}), which fails to meet the fidelity requirements of production anti-fraud systems. Conventional quantization-aware training (QAT) requires reconstructing the full pipeline of supervised fine-tuning (SFT), reinforcement learning (RL), and model merging, increasing both engineering complexity and the risk of destroying the conversational style and refusal boundaries established during alignment. The NVFP4 QAD report from NVIDIA Nemotron (Jan 2026)~\cite{3} suggests a new direction: pure KL-divergence distillation aligning student outputs to a BF16 teacher's distribution can recover 97.0\%--99.4\% of the full-precision performance, providing a methodological foundation for mitigating low-bit quantization degradation in lightweight edge models.

Nevertheless, directly transferring the native QAD paradigm to telecommunication fraud detection faces three challenges:

\textbf{(C1) Generalisation under limited-data conditions:} Anti-fraud conversational corpora are heavily privacy-constrained, restricting the size of high-quality cross-modal training data, and demanding empirical verification of QAD's robustness under limited data coverage.

\textbf{(C2) Hardware incompatibility of multimodal fusion with ultra-low-bit quantization:} Accurate fraud determination requires joint reasoning across text, acoustic, URL, and metadata modalities, yet existing QAD work targets only single-modality (text or vision) models; existing multimodal anti-fraud systems (e.g., SAFE-QAQ) ignore hardware compatibility and feature distribution alignment under NVFP4 and GGUF ultra-low-bit formats.

\textbf{(C3) Stringent on-device latency under hardware--software co-design:} The end-to-end alert latency must remain under 300~ms (median); pure low-bit weight quantization alone cannot break the latency wall imposed by current mobile compute, necessitating joint optimisation with domain-specialised speculative decoding.

To address these challenges, this paper extends the NVFP4 QAD methodology along four dimensions:

\textbf{(1) Pure KL homologous distillation and dual-track quantization.} We adopt pure KL-divergence distillation with a homologous (BF16) teacher, replacing the conventional QAT mixed objective with a single KL loss. A dual-track scheme pairs cloud-side NVFP4 (Blackwell-native) with on-device GGUF-Q4\_K\_M (ARM/Snapdragon-optimised).

\textbf{(2) Output-Variance Freeze (OV-Freeze) regularization.} We constrain the quantized student's $\{q,k,v,o\}$-proj output variance to the BF16 teacher baseline via scalar layerwise rescaling, with EMA variance estimation to suppress mini-batch noise. Appendix~A provides a formal analysis of gradient compatibility and Lipschitz convergence.

\textbf{(3) Irreversible acoustic embedding and multimodal fusion.} We construct $F_v = \mathcal{C}\!\left(f_{\mathrm{mfcc}},\,\psi(W_{\mathrm{proj}}\,\bar h_w)\right) \in \mathbb{R}^{128}$ by concatenating time-averaged MFCC (64-d) with globally average-pooled Whisper-tiny encoder features (64-d). GLO and model-inversion attacks yield WER~$\ge 0.95$; Appendix~A.3 provides an empirical non-invertibility guarantee validated via reconstruction attacks, confirming that $F_v$ is not invertible to intelligible speech under known attacks. L-BFGS-optimised weighted sigmoid fusion of four modalities is validated via user-stratified 5-fold CV.

\textbf{(4) Domain speculative decoding and edge--cloud co-deployment.} We fine-tune a 124M-parameter draft model (Qwen2-0.1B) on anti-fraud corpora, lifting $\alpha$ from 0.78 to 0.86. Measured speedups: $3.49\times$ (H100), $3.32\times$ (SD8G3); end-to-end alert latency P50/P99 = 268/342\,ms. Unified evaluation across TAF-28k, AdvFraud-3k, and ChiFraud yields $F_1 = 0.923$, $0.875$, $0.860$.

The remainder of the paper proceeds as follows. \S~2 reviews related work; \S~3 details the system architecture and methodology; \S~4 reports system- and component-level experiments; \S~5 discusses limitations; \S~6 concludes. Appendix~A provides formal analysis of OV-Freeze and the irreversibility argument for the acoustic embedding.

% ============================================================================
% 2. RELATED WORK
% ============================================================================
\section{Related Work}

\subsection{Quantization-Aware Distillation (QAD)}

The knowledge distillation paradigm introduced by Hinton et al.~\cite{4} transfers knowledge from a high-capacity teacher to a compact student via soft labels. Polino et al.~\cite{5} first integrated KD into weight quantization, founding quantized distillation. Mishra and Marr~\cite{6} demonstrated that low-precision networks can fit the output representations of full-precision teachers to substantially improve accuracy. Liu et al.~\cite{7} proposed LLM-QAT, validating data-free quantization-aware distillation using model-generated data. Du et al.~\cite{8} proposed BitDistiller, combining asymmetric quantization with self-distillation objectives to mitigate sub-4-bit LLM degradation.

The NVFP4 QAD methodology from NVIDIA Nemotron~\cite{3} serves as our primary reference. Compared to conventional QAT, QAD better preserves the base capabilities of densely RL-aligned models. Homologous teacher alignment (BF16 self-distillation) outperforms cross-scale configurations (9B$\rightarrow$9B distillation surpasses 12B$\rightarrow$9B), and KL-divergence proves more effective than MSE for distribution matching. The approach exhibits strong generalisation: distillation on math-only corpora simultaneously recovers code abilities, while random-token inputs do not cause model collapse.

Recent advanced PTQ approaches have made notable progress. Lin et al.~\cite{9} proposed AWQ, protecting critical-channel weights via activation-aware scaling; Frantar et al.~\cite{10}'s GPTQ uses second-order Taylor information for layer-wise weight reconstruction; Liu et al.~\cite{11}'s SpinQuant suppresses activation outliers via learnable orthogonal rotations; Ashkboos et al.~\cite{12}'s QuaRot achieves outlier-free 4-bit inference via Hadamard transforms. While these PTQ methods excel at 8-bit quantization, they exhibit 6.5\%--8.5\% accuracy gaps relative to QAD at 4-bit. This work extends the NVIDIA QAD paradigm to telecommunication-fraud multimodal tasks and introduces the OV-Freeze regularizer to further suppress low-bit quantization noise.

\subsection{Telecommunication Fraud Detection}

Early anti-fraud detection~\cite{13,14,15} relied on keyword-matching rule engines or shallow ML classifiers (SVM, RF). BERT-Fraud~\cite{16} achieved $F_1 = 0.876$ on a Chinese anti-fraud SMS dataset, but lacks on-device deployment support and ignores privacy and acoustic-modality fusion.

Ma et al.~\cite{1} released TeleAntiFraud-28k (TAF-28k), the first open-source audio-text dual-modality dataset with complex reasoning annotations targeting the telecommunication anti-fraud domain. The corpus contains 28{,}511 anonymised voice-text pairs totalling over 307 hours of audio. Real call recordings are first transcribed via automatic speech recognition (ASR), anonymised, and then regenerated via TTS to preserve realistic speech distributions while protecting privacy. The dataset is divided into three core tasks (scenario classification, fraud detection, fraud type classification) with the matching TeleAntiFraud-Bench evaluation protocol.

Wang et al.~\cite{2} proposed SAFE-QAQ, an RL-based end-to-end audio--text multimodal anti-fraud detection framework. It performs end-to-end audio-text joint processing, eliminating cascade errors from ASR transcription, and uses rule-based reasoning rewards to guide layered fraud-pattern recognition, achieving early-warning capability via call-time dynamic risk scoring. According to~\cite{2}, SAFE-QAQ processes approximately $7\times10^4$ calls per day in a reported production deployment (the paper is a preprint and this production figure is cited as reported by the authors without independent verification); however, its high compute cost (7B-scale parameters, server-side inference) limits on-device deployment, which is the work's stated principal limitation.

Our work builds on both directions: TAF-28k provides the data and evaluation protocol; SAFE-QAQ establishes the multimodal-framework foundation; this paper focuses on on-device ultra-low-bit quantization and hardware--software co-deployment. Table~\ref{tab1} summarises the key differences.

\begin{table}[htbp]
\centering\small
\caption{Comparison with related work}
\label{tab1}
\renewcommand{\arraystretch}{1.15}
\setlength{\tabcolsep}{6pt}
\begin{tabular}{>{\centering\arraybackslash}m{2.2cm}>{\raggedright\arraybackslash}m{3.5cm}>{\raggedright\arraybackslash}m{3.5cm}>{\raggedright\arraybackslash}m{5cm}}
\toprule
\textbf{Dimension} & \textbf{TAF-28k~\cite{1}} & \textbf{SAFE-QAQ~\cite{2}} & \textbf{This work (QAD-MultiGuard)} \\
\midrule
Core positioning & Dataset + benchmark & RL end-to-end framework & Edge--cloud deployed system \\
Audio handling & ASR + TTS regeneration & End-to-end audio-text & On-device 128-d $F_v$ (no upload) \\
Quantization & N/A & N/A & NVFP4 + Q4\_K\_M dual-track \\
Deployment & --- & Server-side & Edge--cloud ($<$300 ms) \\
Privacy & Anonymised & Unspecified & $F_v$ non-invertible + DP \\
Training & SFT & RL + slow-think reward & QAD pure-KL distillation \\
Empirical scale & Bench evaluation & 70k calls/day & 5{,}000-user phased deployment \\
\bottomrule
\end{tabular}
\end{table}

\subsection{On-device Large Language Model Inference}

On-device LLM inference is constrained by two factors: model storage size and decoding latency. The GGUF format~\cite{17} compresses 0.5B-scale models to $\sim$240~MB via Q4\_K\_M mixed-precision quantization (4-bit weights with 6-bit critical-layer activations), enabling smooth execution on mobile ARM SoCs. On the server side, vLLM~\cite{18} dramatically improves concurrent throughput via PagedAttention~\cite{19}.

To break the latency bottleneck of autoregressive decoding, Leviathan et al.~\cite{20} proposed speculative decoding, where a lightweight draft model rapidly predicts multiple candidate tokens and the target model verifies them in parallel, achieving 2--3$\times$ end-to-end speedup. Chen et al.~\cite{21} derived the analytic relation between average token acceptance $\alpha$ and sequence speedup:
\begin{equation}
\text{Speedup} = \frac{1-\alpha^{\gamma+1}}{1-\alpha},
\end{equation}
where $\gamma$ is the draft model's single-pass speculation length. Rather than using a generic draft model, we fine-tune on anti-fraud corpora to exploit the stereotyped structure of fraud scripts, lifting $\alpha$ from 0.78 to 0.86.

\subsection{Privacy-Preserving Voice Features}

China's PIPL explicitly mandates that raw biometric information must not leave user terminals without strict authorisation and security assessment. To anonymise speaker identity in local closed-loop processing, Tomashenko et al.~\cite{22}, building on the Voice Privacy Challenge, proposed x-vector substitution for acoustic identity obfuscation. However, Bora et al.~\cite{23}'s Generative Latent Optimization (GLO) attack demonstrated that some traditional anonymisation embeddings remain vulnerable to inverse reconstruction.

This paper proposes a strongly non-invertible acoustic feature construction: time-averaged MFCC and globally average-pooled Whisper-tiny encoder features are concatenated into a 128-d embedding $F_v$; both compression steps destroy frame-level temporal dependencies, yielding WER~$\ge 0.95$ under GLO and model-inversion attacks. Compared to the x-vector substitution of~\cite{22} (whose residual invertibility was demonstrated in~\cite{23}), our method severs both the semantic and speaker-identity channels at the physical compression level. Full details appear in \S~3.3 and Appendix~A.3.

% ============================================================================
% 3. METHOD
% ============================================================================
\section{System Architecture and Method}

\subsection{System Architecture and Threat Model}

\subsubsection{System Architecture}

QAD-MultiGuard adopts a three-tier edge--cloud defence architecture (Figure~\ref{fig1-en}):

\begin{figure}[htbp]
\centering
\includegraphics[width=\linewidth]{Fig/fig1_architecture.png}
\caption{QAD-MultiGuard three-tier edge--cloud collaborative architecture.}
\label{fig1-en}
\vspace{2pt}
\end{figure}

\textbf{Tier-1: On-device detection.} Deployed on Snapdragon~8~Gen~3 and equivalent mobile platforms, this tier runs a lightweight student model (240~MB, Q4\_K\_M quantized), which constructs an SMS feature extractor (12-d), a URL structure-feature extractor (6-d), and a non-invertible acoustic embedding extractor. The acoustic feature vector $\bm{F}_v\in\mathbb{R}^{128}$ ensures that raw audio and SMS plain text remain on-device, fulfilling PIPL requirements.

\medskip
\noindent\textbf{Data-flow boundary.} We delineate precisely what leaves the device. For voice calls, a Whisper-tiny encoder runs on-device \emph{without its decoder head}---no ASR transcript is produced. Its encoder output is pooled and projected into the 64-d Whisper branch of $\bm{F}_v$. For SMS, the on-device student LLM processes 12-d structural features and outputs only a scalar $r_{\text{text}}\in[0,1]$. URL and metadata risks are likewise on-device scalars. The cloud Tier-2 therefore receives: (i) the 128-d $\bm{F}_v$ (non-invertible, Proposition~A.3); (ii) four scalar risk scores; and (iii) intermediate hidden states for CoT cross-validation. Raw audio, ASR transcripts, and SMS plaintext never leave the device.

\textbf{Tier-2: Cloud deep reasoning.} Deployed on the NVIDIA Blackwell GPU cluster, running an NVFP4-quantized 0.5B LLM cluster for high-concurrency chain-of-thought (CoT) reasoning and multimodal cross-validation. The cloud and on-device models share the same backbone architecture (Table~\ref{tab2-en}), constituting a fully homologous edge--cloud architecture: the cloud model undergoes NVFP4 quantization while the edge model uses Q4\_K\_M quantization, both distilled from the same BF16 teacher. Although 0.5B-scale models are generally considered insufficient for open-domain CoT reasoning, telecommunication fraud detection operates within a sharply constrained semantic domain: (i) the top-50 high-frequency fraud terms cover 73.4\% of fraud-text samples (\S~3.4), limiting the effective reasoning vocabulary to a few hundred domain-specific tokens; (ii) fraud scripts follow highly stereotyped templates (e.g., ``impersonate authority'' $\rightarrow$ ``create urgency'' $\rightarrow$ ``request transfer''), reducing the CoT reasoning depth to 2--4 steps. Critically, the cloud model does not directly receive raw text; it operates on the guided hidden states transmitted from the edge, which fuse textual and acoustic semantic information into a modality-aligned representation via the homologous architecture. The internal CoT reasoning is therefore performed on these semantic patterns, not on plaintext, and is semantically equivalent to the following interpretation: ``[Interpreted Semantic Pattern] Authority-impersonation template (``police officer'') + urgency coercion (``immediate transfer'') + target-phrase match (``safe account'' $\rightarrow$ template \#A07).
[Reasoning] (1) Authority impersonation via semantic pattern matching. (2) Urgency coercion via imperative construction detection. (3) Target phrase aligns with known fraud template \#A07. $\rightarrow$ HIGH RISK.'' Table~\ref{tab:cot-ablation-en} confirms that removing CoT reduces $F_1$ by 1.8 points (0.923$\rightarrow$0.905).

\textbf{Tier-3: Multimodal risk fusion.} Also running on-device, this tier optimises the decision boundary off-line via L-BFGS and applies weighted sigmoid fusion to combine the four modality risk scores, producing a real-time tri-level risk verdict (Safe / Medium / High) that triggers on-device UI alerts.

\subsubsection{Formalised Threat Model}

We formalise the threat scenario as a 4-tuple $\langle \mathcal{A}, \mathcal{K}, \mathcal{G}, \mathcal{C} \rangle$ covering attacker capabilities, prior knowledge, malicious intent, and defence-system safety constraints.

\textbf{(1) Attacker capability ($\mathcal{A}$):} The attacker is positioned at an untrusted network boundary (compromised access point or proxy gateway), acting as an active or passive man-in-the-middle (MITM). They can intercept and parse the acoustic embedding $\bm{F}_v$ and associated textual features transmitted over the edge--cloud channel, and query the cloud inference API in black-box fashion. They cannot, however, breach the on-device trusted execution environment (TEE).

\textbf{(2) Attacker knowledge ($\mathcal{K}$):} We assume the attacker has full white-box knowledge of on-device feature extractors (including the projection matrix $\bm{W}_{\mathrm{proj}}$) and has access to historical $\bm{F}_v$ sequences from past interception; in contrast, the attacker has only black-box knowledge of the cloud LLM's topology, weights, and hyperparameters.

\textbf{(3) Attacker goals ($\mathcal{G}$):} Three heterogeneous threat paths are considered.
\begin{itemize}
    \item \textbf{$\mathcal{G}_1$ (privacy reconstruction):} Train an inverse mapping to recover raw waveform or high-fidelity transcription from intercepted $\bm{F}_v$.
    \item \textbf{$\mathcal{G}_2$ (adversarial evasion):} Inject semantic-level adversarial prompts or context variants to mislead the cloud LLM's risk-decision pipeline.
    \item \textbf{$\mathcal{G}_3$ (identity fraud):} Replay intercepted feature vectors or forge credentials to induce false authorisation at the cloud.
\end{itemize}

\textbf{(4) Safety constraints ($\mathcal{C}$):}
\begin{itemize}
    \item Against $\mathcal{G}_1$: reverse-reconstructed speech must maintain $\text{WER}\ge 0.90$, and feature extraction must satisfy $\epsilon$-local differential privacy ($\epsilon$-LDP, with $\epsilon = 1.5$).
    \item Against $\mathcal{G}_2$: $F_1$ degradation on the AdvFraud-3k adversarial set must remain $\le 6$ percentage points.
    \item Against $\mathcal{G}_3$: cross-modal validation must use call metadata and other heterogeneous features.
\end{itemize}

\textbf{(5) Scope exclusion.} This threat model focuses on attackers operating fully automated generative-fraud pipelines ($\text{ASR}\!\rightarrow\!\text{LLM}\!\rightarrow\!\text{TTS}$). We exclude physical-signal-layer adversarial perturbations (e.g., ASRJam~\cite{19}) because (i) multimodal fusion provides inherent robustness to single-modality perturbation and (ii) AdvFraud-3k (\S~4.3) covers 8 semantic-level adversarial variants. Physical-layer robustness is left to future work (\S~5.3).

\subsection{Quantization-Aware Distillation (QAD)}

\subsubsection{Pure KL-divergence loss}

To best preserve the prior distribution acquired across multi-stage alignment during low-bit quantization, QAD-MultiGuard introduces a pure KL-divergence-driven distillation paradigm. Given input sequence $x$ and vocabulary $\mathcal{V}$, let $p_{\text{teacher}}(y|x)$ denote the BF16 teacher's output and $p_{\text{student}}(y|x)$ the quantized student's output. The QAD loss is
\begin{equation}
L_{\mathrm{QAD}} = D_{\mathrm{KL}}\!\bigl( p_{\text{teacher}}(y|x)\,\|\,p_{\text{student}}(y|x) \bigr) = \sum_{y\in\mathcal{V}} p_{\text{teacher}}(y|x)\log\!\frac{p_{\text{teacher}}(y|x)}{p_{\text{student}}(y|x)}.
\end{equation}
The Softmax temperature is uniformly fixed to $T = 1$ throughout distillation, enabling unbiased approximation of the teacher's original probability stream.

Conventional QAT~\cite{24,25,26} typically retains the SFT-stage cross-entropy, fitting task targets under discretization constraints. For models progressively shaped by ``SFT fraud-text alignment + RL safety alignment'', QAT induces severe distribution drift: cross-entropy cannot constrain non-top-1 token probabilities, leading to catastrophic degradation of the RL-stage safety refusal boundary and reply style. Empirically, conventional QAT produces a KL divergence of 0.311 against the BF16 teacher, whereas pure-KL QAD compresses this to 0.004 while inheriting teacher safety behaviour.

\subsubsection{Homologous teacher self-distillation}

Conventional KD follows the ``large teacher $\rightarrow$ small student'' cross-scale capability transfer. For low-bit quantization of large models, however, heterogeneous teachers are typically inferior to homologous self-distillation, because QAD's core objective is to repair output-distribution perturbations from weight discretization---not to inject new task knowledge. We therefore use the native Qwen2.5-0.5B-Instruct (BF16)~\cite{27,28,29} as the homologous teacher.

\subsubsection{Edge--cloud co-quantization}

Rather than relying on a single-track cloud-only quantization paradigm (e.g., NVFP4 operators on Blackwell), we adopt an edge--cloud dual-track scheme that couples the server cluster's high-performance compute with mobile hardware's low-level instruction sets. The deployment parameters are summarised in Table~\ref{tab2-en}.

\begin{table}[h]
\centering
\caption{Edge--cloud dual-track quantization parameters}
\label{tab2-en}
\begin{tabular}{l>{\raggedright\arraybackslash}p{6.5cm}>{\raggedright\arraybackslash}p{4.7cm}}
\toprule
\textbf{Component} & \textbf{Server (NVFP4)} & \textbf{On-device (Q4\_K\_M)} \\
\midrule
Backbone & Qwen2.5-0.5B-Instruct & Qwen2.5-0.5B-Instruct \\
Parameters & 494M (FP16) & 494M (FP16) \\
Teacher & BF16 self-distillation & BF16 self-distillation \\
Quant format & NVFP4 (block=16, E4M3 scale) & Q4\_K\_M (4/6-bit GGUF) \\
Model size & 248 MB & 240 MB \\
Platform & Blackwell GPU + vLLM & ARM / Snapdragon \\
Speedup vs FP8 & 2.1$\times$ & --- \\
Recovery vs BF16 & \textbf{99.1\%} & \textbf{98.5\%} \\
\bottomrule
\end{tabular}
\end{table}

The cloud cluster adopts NVFP4, integrating three design elements: 16-element micro-block scaling ($\text{block\_size}=16$), FP8 E4M3 scale storage, and coordinated tensor-level FP32 global scaling. Micro-block scaling maps local dynamic range tightly into the FP4 representation range, avoiding saturation from long-tail distributions and converging quantization noise; FP8 E4M3 provides high-resolution scale representation, mitigating truncation/rounding bias in ultra-low-bit space; the FP32 global scale remaps tensor distribution at the macro level, guiding micro-block scales to cover the effective dynamic range of activations and weights.

On the mobile side, we adopt the GGUF ecosystem's Q4\_K\_M grouped quantization: predominantly 4-bit compact representation with multi-dimensional group-scale hyperparameters controlling precision loss; operator paths bind ARM SIMD vectorised instruction sets for operator-level high-throughput parallelism.

\medskip
\noindent\textbf{Model size note.} The 8\,MB difference between the two tracks (NVFP4: 248\,MB, Q4\_K\_M: 240\,MB) is attributable to differing scale metadata overheads. NVFP4 uses fine-grained per-16-element micro-block scaling with FP8 E4M3 plus a per-tensor FP32 scale, whose metadata adds $\sim$8\,MB relative to GGUF's coarser super-block structure. This 3.3\% overhead is the price of NVFP4's finer-grained dynamic range adaptation, which contributes to its slightly higher recovery rate (99.1\% vs.\ 98.5\%).

\subsubsection{Numerical Behaviour Emulation (NBE) Protocol}

We adopt a Numerical Behaviour Emulation (NBE) protocol on H100 to build equivalent closed-loop alignment to the cloud NVFP4 pipeline, decoupling numerical-precision verification from hardware-specific performance gains. Fake-quantization (QDQ) operators are embedded in the forward graph: weights and activations are computed in FP16/BF16, but pass through a ``quantize--dequantize'' step before each linear layer, faithfully reproducing NVFP4 rounding and saturation error at FP16 precision. This scheme is behaviour-aligned with NVIDIA TensorRT-LLM (v0.13)~\cite{NVIDIA2024TRTLLMReleaseNotes,NVIDIA2024TensorRTLLMDocs}.

Concretely, each weight tensor $\bm{W}$ undergoes
\begin{equation}
\widehat{\bm{W}} = \mathrm{clamp}\!\bigl(\mathrm{round}(\bm{W}/s),\,q_{\min},\,q_{\max}\bigr)\cdot s,
\end{equation}
with dual-layer scale $s = s_{\text{block}}\!\cdot\!s_{\text{tensor}}$: $s_{\text{block}}$ (block=16, FP8-E4M3) and $s_{\text{tensor}}$ (FP32 global). The dequantized weights are injected into matmul under FP16/BF16, exposing low-bit error propagation. The protocol explicitly excludes Blackwell-specific hardware effects (2:4 sparsity, tcgen rounding, memory hierarchy); all throughput and latency are measured on H100.

Functional correctness is established via two-step calibration: (1) cross-comparison against the NVIDIA NVFP4 reference checkpoint (Nemotron Nano 9B V2) on MMLU~\cite{30}, GSM8K~\cite{31}, and HumanEval~\cite{32} yields absolute accuracy deviation $<0.3\%$, verifying QDQ-chain correctness; (2) 5-seed repetition on our 0.5B student yields seed-wise $F_1$ $\sigma = \pm 0.006$ on TAF-28k. A scale gap is present (9B reference vs.\ 0.5B student), but the 0.3\% deviation validates only QDQ-chain correctness, and since QDQ operators are \emph{model-architecture-agnostic} (operating identically on any dense linear layer), the validated chain transfers to 0.5B without modification. The $\pm 0.006$ $F_1$ standard deviation bounds any scale-transfer residual. All results are therefore reported as mean $\pm$ std.\ across 5 seeds.

\subsubsection{Output-Variance Freeze (OV-Freeze) Regularization}

Low-bit discretization causes severe output-distribution drift on the four sensitive projection layers of self-attention ($\ell \in \{q, k, v, o\}_{\mathrm{proj}}$), with empirically observed output-variance relative increase up to $+18.2\%$. This variance shift disrupts attention-score scale balance, degrading deep-network numerical stability. OV-Freeze, a lightweight regularization module, is introduced in the post-training phase to constrain the quantized student's sensitive-layer output variance to the homologous BF16 teacher's statistical level.

OV-Freeze activates only in the final 30\% of post-training optimisation steps. It depends on offline-calibrated teacher (BF16) output-variance constants $\sigma^2_{\mathrm{BF16},\ell}$. During forward pass, the student's layer output $\bm{y}_\ell$ undergoes adaptive variance recalibration:
\begin{equation}
\label{eq:ovf-rescale-en}
\bm{y}'_\ell = \bm{y}_\ell\cdot c_\ell, \qquad c_\ell = \sqrt{\frac{\sigma^2_{\mathrm{BF16},\ell}}{\mathrm{Var}(\bm{y}_\ell) + \epsilon}}.
\end{equation}
where $\epsilon = 10^{-6}$ is a numerical-stability constant. This operation is equivalent to applying layer-wise scalar rescaling without introducing new structural assumptions or altering attention topology.

\textbf{(1) Gradient compatibility and implementation.} In the actual forward--backward graph, the variance branch is detached during backpropagation, so $c_\ell$ acts as a constant (input-independent) scaling factor and the gradient simplifies to:
\begin{equation}
\frac{\partial L_{\mathrm{QAD}}}{\partial\bm{y}_\ell} = c_\ell\cdot\frac{\partial L_{\mathrm{QAD}}}{\partial\bm{y}'_\ell}.
\label{eq:ovf-detach-en}
\end{equation}
This design preserves the original KL-gradient direction exactly, introducing only per-layer magnitude-level scaling. Although the full gradient (without detachment) would contain an additional Jacobian term $\partial c_\ell/\partial\bm{y}_\ell$,
\begin{equation}
\frac{\partial L_{\mathrm{QAD}}}{\partial\bm{y}_\ell} = c_\ell\cdot\frac{\partial L_{\mathrm{QAD}}}{\partial\bm{y}'_\ell} + \left\langle\frac{\partial L_{\mathrm{QAD}}}{\partial\bm{y}'_\ell},\,\bm{y}_\ell\right\rangle\cdot\frac{\partial c_\ell}{\partial\bm{y}_\ell},
\label{eq:ovf-full-grad-en}
\end{equation}
the detached form~\eqref{eq:ovf-detach-en} is the operational forward--backward path. The full Jacobian~\eqref{eq:ovf-full-grad-en} is analysed theoretically in Proposition~A.1 (Appendix~A.1) as a worst-case boundedness guarantee, confirming that even without detachment the correction factor cannot cause gradient explosion; the practical detached version is a strict special case that inherits this guarantee.

\textbf{(2) Batch noise suppression via EMA.} Under small batch ($\text{batch\_size}=8$), directly using batch variance $\mathrm{Var}_{\mathrm{batch}}(\bm{y}_\ell)$ causes severe early-stage oscillation (coefficient of variation CV~$\approx 0.43$). We introduce exponential moving average (EMA) for stable long-window variance estimation:
\begin{equation}
\mathrm{Var}^{(t)}_{\mathrm{EMA}}(\bm{y}_\ell) = \rho\cdot\mathrm{Var}^{(t-1)}_{\mathrm{EMA}}(\bm{y}_\ell) + (1-\rho)\cdot\mathrm{Var}^{(t)}_{\mathrm{batch}}(\bm{y}_\ell), \qquad \rho = 0.95.
\end{equation}
The EMA buffer is initialised to the teacher baseline $\sigma^2_{\mathrm{BF16},\ell}$. With $\rho = 0.95$, the effective half-life is $\sim$14 steps ($\sim$110 independent samples), suppressing the variance-estimate CV below 0.05. Forward computation uses $\mathrm{Var}_{\mathrm{EMA}}$; the EMA buffer is detached in backprop.

\textbf{(3) Ablation.} Batch-level estimation induces $+1.4$ PPL fluctuation in the first 500 steps (final $F_1 = 0.918$); global static statistics lack adaptive tracking of teacher--student drift ($F_1 = 0.919$); EMA achieves the best trade-off ($F_1 = 0.923$, PPL fluctuation $+0.18$; see Table~\ref{tab:ovf_var_ablation-en}). Across both NVFP4 and Q4\_K\_M tracks, OV-Freeze contributes $+0.7$ and $+0.6$ $F_1$ percentage points over the un-activated baselines ($\Delta\text{PPL}=-0.11$; see \S~4.5).

\subsection{Privacy-Preserving Acoustic Embedding}

The system observes PIPL data-compliance constraints. The decentralised topology extracts low-dimensional, non-invertible acoustic features locally, blocking original-waveform exfiltration while feeding downstream risk classification. The privacy-preserving acoustic embedding is defined as
\begin{equation}
\bm{F}_v = \mathcal{C}\!\left(\bm{f}_{\mathrm{mfcc}},\,\psi(\bm{W}_{\mathrm{proj}}\bar{\bm{h}}_w)\right)\in\mathbb{R}^{128}.
\end{equation}
Here $\mathcal{C}(\cdot)$ denotes concatenation and $\psi(\cdot)$ denotes vector-wise normalization. The MFCC vector $\bm{f}_{\mathrm{mfcc}}\in\mathbb{R}^{64}$ is extracted at 16 kHz sampling with 64-band mel filter banks, 10 ms hop, 25 ms window. The sequence pooled representation $\bar{\bm{h}}_w\in\mathbb{R}^{384}$ is derived from Whisper-tiny encoder global average pooling, mapped to a 64-d subspace via the learnable matrix $\bm{W}_{\mathrm{proj}}\in\mathbb{R}^{64\times 384}$. This projection is jointly optimised with the student model in the QAD fine-tuning phase, with row-wise normalisation preventing projection collapse.

\subsubsection{Streaming sliding-window implementation}

Real-time alert latency is achieved via streaming sliding windows. The on-device ring buffer maintains a $W = 3$\,s temporal audio window, advancing in steps $\Delta = 1$\,s (66.7\% inter-window overlap). Each full window triggers $\bm{F}_v$ extraction and a single Tier-2 risk-score iteration (on-device student inference $\rightarrow$ feature upload $\rightarrow$ cloud NVFP4 CoT inference $\rightarrow$ score download $\rightarrow$ on-device fusion). The per-window incremental inference latency has measured P50 $=268$~ms ($<500$~ms incremental budget, matching latency notion (i) in \S~1.1); the first-alert latency is $W + t_{\mathrm{infer}}\approx 3.3$~s (latency notion (ii)), after which risk scores refresh at a 1\,s cadence.

\medskip
\noindent\textbf{Latency breakdown.} The 268\,ms P50 decomposes into: (a) on-device feature extraction ($\bm{F}_v$ + text/URL/meta features, $\sim$12\,ms); (b) encrypted feature upload over Wi-Fi~6 ($\sim$5\,ms RTT, 128-d $F_v$ + 4 scalars = $\sim$1\,KB payload); (c) cloud NVFP4 CoT inference with domain-adapted speculative decoding ($\sim$230\,ms, Table~\ref{tab:speculative_decoding-en} network-excluded single-query latency); (d) score download and on-device fusion ($\sim$21\,ms). Measurements are under single-user, single-call load on a production-grade 5G/Wi-Fi~6 network environment (lab deployment). Under concurrent multi-call load (simulated 50 simultaneous calls), P50 degrades to $\sim$310\,ms (P99 $\sim$410\,ms) due to cloud-side queueing, remaining within the 500\,ms incremental budget. All latency figures exclude the initial 3\,s acoustic window accumulation (which is a physical constraint, not a processing delay).

\medskip
\noindent Critically, the ``time averaging'' in this work is defined \emph{within a single sliding window} ($W=3$~s, $\sim$300 MFCC frames), \emph{not} as a cumulative average over the entire call. Once the first full window is built ($\approx 3.3$~s), streaming alerts fire without waiting for the call to end, resolving the apparent tension between real-time constraints and time-averaged acoustic features.

The window span $W$ governs a Pareto trade-off between alert latency, MFCC statistical robustness, and temporal resolution. Ablation over $W\in\{1,2,3,5,10\}$\,s shows that $W = 3$\,s is the empirical Pareto optimum: $W = 1$\,s amplifies high-frequency statistical noise from sparse frame counts, dropping $F_1$ to 0.897; $W = 5$\,s offers only marginal accuracy gain ($F_1 = 0.924$) at the cost of latency; $W = 10$\,s severely degrades real-time responsiveness.

Inter-window score smoothing uses exponential moving average:
\begin{equation}
r_t = \beta r_{t-1} + (1-\beta)r_t^{(\mathrm{new})}, \qquad \beta = 0.4.
\end{equation}
High-risk interception is triggered when two consecutive aggregated scores satisfy $r\ge 0.70$.

\subsubsection{Non-invertibility argument and differential privacy extension}

The construction of $\bm{F}_v$ explicitly strips frame-level high-frequency temporal features, severing the inverse-recovery path. This non-invertibility arises from dual temporal compression: (i) MFCC global temporal averaging ($\sim$300 MFCC frames averaged into a single 64-d vector) compresses phoneme-level temporal structure into static statistical moments; (ii) Whisper encoder output undergoes sequence-level average pooling followed by a rank-limited ($\text{rank}\le 64$) projection, eliminating token-level fine-grained context. We emphasise that this is an \emph{engineering} non-invertibility guarantee rather than a formal information-theoretic claim---because $F_v$ retains task-relevant acoustic statistics (e.g., energy variance, pitch range), a small positive mutual information $I(x;F_v)>0$ necessarily exists. The critical question is whether this residual information suffices for any practical inverse attack, which we assess empirically below.

Embedding recoverability is quantitatively bounded via white-box GLO attacks and black-box model inversion. On a 100-segment test set of real call recordings (3\,s each), white-box and black-box inverse attacks yield WER of 0.95 and 0.97 respectively; speaker-identification accuracy drops to 8.3\%, approaching the 10\% random baseline. These results demonstrate that, despite $I(x;F_v)>0$ theoretically, the residual information is insufficient for any known inverse attack to recover intelligible speech or identifiable speaker characteristics, satisfying PIPL compliance requirements for practical deployment. The privacy quantitative metrics between $W = 1$\,s and $W = 5$\,s differ by less than 0.5\%, confirming the orthogonality of window scale to the non-invertibility properties.

For extreme-privacy scenarios, the system supports adaptive Gaussian noise injection to construct $(\epsilon, \delta)$-DP variants:
\begin{equation}
\bm{F}_v' = \bm{F}_v + \mathcal{N}(0, \sigma^2).
\end{equation}
By controlling $\ell_2$ sensitivity and the Gaussian scale $\sigma$, the system can strictly converge to a preset $(\epsilon,\delta)$-DP bound.

\subsection{Domain-Adapted Speculative Decoding}

Anti-fraud text exhibits high pattern-density. According to Ma et al.'s statistics on TAF-28k, the top-50 high-frequency fraud terms cover 73.4\% of fraud-text samples. Building on this prior knowledge, we fine-tune a 124M-parameter draft model (Qwen2-0.1B) on 50K anti-fraud streaming utterances to align with the cloud target model.

In speculative verification, the expected token yield per iteration is
\begin{equation}
\mathbb{E}[N] = \frac{1-\alpha^{\gamma+1}}{1-\alpha},
\end{equation}
where $\gamma$ is the draft model's single-pass speculation length. At $\gamma = 5$, the generic draft model has $\alpha\approx 0.78$ and expected yield 3.52 tokens; domain-adapted fine-tuning lifts $\alpha$ to 0.86 and expected yield to 4.25. This theoretical speedup is partially offset by KV-cache access and parallel verification overheads in the hardware-measured implementation (\S~4.8): measured H100 speedup is $3.49\times$, SD8G3 is $3.32\times$ (Table~\ref{tab:speculative_decoding-en}).

Sensitivity analysis over $\gamma\in\{3,5,7,10\}$ in \S~4.9 establishes $\gamma = 5$ as the Pareto optimum between text-generation efficiency and KV-cache memory. Relative to $\gamma = 5$: $\gamma = 7$ yields only +17\% expected token throughput but 40\% higher KV-cache memory; $\gamma = 3$ reduces memory but loses 24\% throughput.

\subsection{Multimodal Risk Fusion}

\subsubsection{Modality risk scoring}

Final risk determination relies on decision-level fusion of text, acoustic, URL, and metadata. All modality risk inputs are mapped to $[0, 1]$:

\textbf{Text risk $r_{\mathrm{text}}\in[0,1]$:} the on-device student LLM produces a sigmoid posterior probability from 12-d SMS structural features and anonymised text summaries.

\textbf{Acoustic risk $r_{\mathrm{audio}}\in[0,1]$:} four acoustic proxy statistics derived from $\bm{F}_v$ (energy variance, tone proxy, urgency proxy, pitch range) are linearly aggregated:
\begin{equation}
r_{\mathrm{audio}} = \alpha_1\cdot\text{energy\_var} + \alpha_2\cdot\text{tone\_proxy} + \alpha_3\cdot\text{urgency\_proxy} + \alpha_4\cdot\text{pitch\_range}.
\end{equation}
with $\bm{\alpha} = [0.35, 0.28, 0.25, 0.12]^T$ learned via $L_2$-regularized logistic regression ($\lambda = 0.01$) on Z-score-standardised features, normalised to $\sum_{i=1}^4\alpha_i = 1$. The two-level weight-learning design---intra-modality $\bm{\alpha}$ fit offline by lightweight logistic regression, inter-modality $\mathbf{w}$ optimised by L-BFGS (\S~3.5.2)---decouples interpretability from cross-modal global optimality.

\textbf{URL risk $r_{\mathrm{url}}\in[0,1]$:} weighted scoring of 6-d URL structural features (domain length, path depth, IP direct access, non-standard port, entropy, short-URL flag).

\textbf{Metadata risk $r_{\mathrm{meta}}\in[0,1]$:} gradient-boosting machine (GBM) classifier on 12-d call metadata (historical report frequency, confirmed-fraud frequency, caller credibility).

\subsubsection{L-BFGS decision fusion}

The global risk score is
\begin{equation}
r = \sigma\bigl(w_{\mathrm{text}}\cdot r_{\mathrm{text}} + w_{\mathrm{audio}}\cdot r_{\mathrm{audio}} + w_{\mathrm{url}}\cdot r_{\mathrm{url}} + w_{\mathrm{meta}}\cdot r_{\mathrm{meta}} + b\bigr),
\end{equation}
where $\sigma(\cdot)$ is the sigmoid function. The parameters $\mathbf{w}$ and $b$ are optimised via L-BFGS on a labelled validation subset.

To assess decision-weight generalisation stability, we perform \emph{user-stratified} 5-fold cross-validation (all calls from one user appear in the same fold, preventing cross-fold leakage). The cross-fold weight distribution is $\mathbf{w} = [0.40 \pm 0.018, 0.30 \pm 0.015, 0.20 \pm 0.012, 0.10 \pm 0.009]^T$. The risk thresholds are: $r\ge 0.70 \Rightarrow$ high risk (interception); $0.35\le r < 0.70 \Rightarrow$ medium risk; otherwise safe.

\subsubsection{Comparison with deep fusion baselines}

To assess the performance ceiling of linear fusion, we introduce a multimodal Transformer baseline: four modality features are projected to a shared latent space, concatenated as four tokens, fed to a 2-layer Transformer encoder, pooled, and classified. Table~\ref{tab:fusion_ablation-en} shows the Transformer yields only $+0.004$ $F_1$ gain ($0.927$ vs.\ $0.923$) at $\sim$42$\times$ parameter inflation and $+7.2$~ms latency; Softmax fusion (mutually-exclusive prior across modalities) degrades $F_1$ by $0.014$. Under edge constraints, our linear Sigmoid fusion approaches the deep-fusion ceiling at negligible cost.

\begin{table}[htbp]
\centering
\caption{Multimodal fusion strategy ablation (TAF-28k test set)}
\label{tab:fusion_ablation-en}
\small
\renewcommand{\arraystretch}{1.2}
\begin{tabular}{l>{\centering\arraybackslash}p{0.12\linewidth}>{\centering\arraybackslash}p{0.16\linewidth}>{\centering\arraybackslash}p{0.15\linewidth}>{\centering\arraybackslash}p{0.13\linewidth}}
\toprule
\textbf{Fusion strategy} & \textbf{$F_1$} & \textbf{Fusion params} & \textbf{Edge latency} & \textbf{Memory} \\
\midrule
Softmax linear & $0.909$ & 5 scalars & $\sim 1$~ms & negligible \\
\rowcolor[gray]{0.93}\textbf{Sigmoid linear (ours)} & $\mathbf{0.923}$ & \textbf{5 scalars} & $\mathbf{\sim 1}$~\textbf{ms} & \textbf{negligible} \\
Multimodal Transformer (2-layer) & $0.927$ & $1.84$~M & $\sim 8.2$~ms & $+12$~MB \\
\bottomrule
\end{tabular}
\end{table}

% ============================================================================
% 4. EXPERIMENTS
% ============================================================================
\section{Experimental Evaluation}

\subsection{Experimental Setup}

\subsubsection{Datasets}

\textbf{TAF-28k~\cite{1}:} 28{,}511 voice-text pairs ($\sim$307 h audio), split 8:1:1 train/val/test. Evaluation follows the TeleAntiFraud-Bench protocol on scenario classification, fraud detection, and fraud type classification.

\textbf{AdvFraud-3k:} 1{,}000 TAF-28k fraud-text samples manually rewritten under 8 adversarial strategies (synonym substitution, syntactic reordering, dialect shift, metaphor) plus 2{,}000 expert-composed novel fraud scripts. Three independent annotators with a fourth adjudicator (Cohen's $\kappa = 0.87$). Used exclusively for evaluation; no train-test leakage.

\textbf{ChiFraud~\cite{chifraud}:} A long-period Chinese network-fraud text benchmark, used as the external long-period out-of-distribution dataset for testing cross-domain generalisation.

\subsubsection{Baselines}

We benchmark against 10 representative low-bit quantization and anti-fraud methods. PTQ methods: AWQ, GPTQ, SpinQuant, QuaRot (all under NVFP4 numerical configuration, \emph{without} knowledge distillation or QAD). Self-distillation: BitDistiller. QAT: standard cross-entropy NVFP4 QAT. Q4\_K\_M PTQ (also without distillation). Domain-specific: BERT-Fraud and SAFE-QAQ.

\medskip
\noindent\textbf{Fairness of PTQ comparisons.} We acknowledge that PTQ methods (AWQ, GPTQ, SpinQuant, QuaRot) are post-training quantization techniques rather than distillation-aware methods, and their comparison against QAD captures the benefit of \emph{distillation-aware} quantization rather than an algorithmic advantage in weight rounding alone. To provide a more controlled reference, we additionally report a ``NVFP4 PTQ + temperature scaling'' variant, where the PTQ model's Softmax temperature is calibrated on a 100-sample held-out set to minimise KL divergence against the BF16 teacher. This simple post-hoc calibration improves PTQ $F_1$ from 0.838 to 0.852 (Table~\ref{tab3-en} footnote), but remains 7.1 points below QAD+OVF, confirming that the primary accuracy recovery stems from distillation-driven distribution alignment rather than temperature rescaling.

\medskip
\noindent\textbf{SAFE-QAQ configuration.} The SAFE-QAQ baseline uses its original configuration (7B-scale parameters, BF16 precision, server-side inference, without quantization) as reported in~\cite{2}. We do not adapt SAFE-QAQ to 0.5B or apply quantization because its RL-based multi-stage training pipeline is not designed for low-bit deployment and would require substantial architectural modification; the comparison therefore evaluates the trade-off between our compact quantized framework and a larger unquantized production-grade system.

\subsubsection{Hardware platforms}

\textbf{Server:} NVIDIA H100 80GB SXM5; NVFP4 numerical emulation via TensorRT-LLM, cross-calibrated against NVIDIA's official NVFP4 reference checkpoints.

\textbf{Edge:} Snapdragon~8~Gen~3 (Hexagon NPU, ARM Cortex-X4, 12\,GB RAM).

\textbf{Training cluster:} $8\times$NVIDIA H100, PyTorch 2.5 with DeepSpeed ZeRO-3. All experiments use 5 random seeds; mean $\pm$ standard deviation reported.

\subsubsection{Training configuration}

QAD hyperparameters follow NVIDIA's official optimiser configuration: cosine LR decay with initial $1\times 10^{-5}$, warm-up first 100 steps, batch size 8, context length 4096, total 2{,}000 steps consuming $\sim$0.5B tokens (1.7\% of the original SFT corpus). OV-Freeze activates in the final 600 steps (30\%) with regularization coefficient $\lambda = 1.0$.

\subsubsection{Evaluation metrics}

\textbf{Recovery rate} $= F_{1,\mathrm{quant}}/F_{1,\mathrm{BF16}}\times 100\%$. \textbf{Classification:} $F_1$, precision, recall, FPR. \textbf{Distribution alignment:} KL divergence between quantized student and BF16 teacher posteriors. \textbf{Latency:} P50 and P99 end-to-end alert latency. \textbf{Privacy strength:} reconstructed-speech WER under white-box GLO and black-box model-inversion attacks.

\medskip
\noindent\textbf{Statistical significance.} All reported $F_1$ values are mean $\pm$ standard deviation over 5 random seeds. BERT-Fraud is deterministic (single seed) and is reported without $\pm$ in its original publication. For key pairwise comparisons (e.g., QAD+OVF vs.\ QAD, OVF activation ratios), we additionally perform a two-sided paired bootstrap test (10{,}000 resamples) on the test-set predictions; results marked as significant satisfy $p<0.01$ unless otherwise noted. The OV-Freeze gain ($+0.7$ $F_1$ points) is significant at $p<0.01$ (bootstrap 95\% CI [$+0.31$, $+1.08$]).

\subsection{Main results}

Table~\ref{tab3-en} presents results on TAF-28k. NVFP4 numerical data is generated on H100; consistency with NVIDIA's official NVFP4 architecture (Nemotron-Nano-9B-V2 as reference) is verified with classification-boundary deviation $\le 0.3\%$.

\begin{table}[h]
\centering
\caption{TAF-28k test-set multimodal fraud detection results}
\label{tab3-en}
\small
\renewcommand{\arraystretch}{1.2}
\begin{tabularx}{\textwidth}{Xccccc}
\toprule
\textbf{Method} & \textbf{$F_1$} & \textbf{Precision} & \textbf{Recall} & \textbf{FPR} & \textbf{Recovery (\%)} \\
\midrule
BF16 (upper bound) & $0.931 \pm 0.005$ & $0.928$ & $0.934$ & $1.6\%$ & $100.0$ \\
\midrule
\multicolumn{6}{l}{\emph{\textbf{Advanced post-training quantization (PTQ)}}} \\
NVFP4 PTQ (max) & $0.838 \pm 0.011$ & $0.847$ & $0.829$ & $2.8\%$ & $90.0$ \\
NVFP4 + AWQ & $0.838 \pm 0.010$ & $0.846$ & $0.830$ & $2.7\%$ & $90.0$ \\
NVFP4 + GPTQ & $0.840 \pm 0.010$ & $0.848$ & $0.832$ & $2.7\%$ & $90.2$ \\
NVFP4 + SpinQuant & $0.838 \pm 0.011$ & $0.846$ & $0.830$ & $2.7\%$ & $90.0$ \\
NVFP4 + QuaRot & $0.838 \pm 0.011$ & $0.846$ & $0.830$ & $2.7\%$ & $90.0$ \\
\midrule
\multicolumn{6}{l}{\emph{\textbf{Self-distillation low-bit quantization}}} \\
NVFP4 + BitDistiller & $0.858 \pm 0.009$ & $0.866$ & $0.850$ & $2.5\%$ & $92.2$ \\
\midrule
\multicolumn{6}{l}{\emph{\textbf{QAT/QAD (ours)}}} \\
NVFP4 QAT & $0.844 \pm 0.014$ & $0.853$ & $0.835$ & $3.1\%$ & $90.7$ \\
NVFP4 QAD & $0.916 \pm 0.007$ & $0.918$ & $0.914$ & $1.9\%$ & $98.4$ \\
\rowcolor[gray]{0.93} \textbf{NVFP4 QAD + OV-Freeze} & $\mathbf{0.923 \pm 0.006}$ & $\mathbf{0.925}$ & $\mathbf{0.921}$ & $\mathbf{1.8\%}$ & $\mathbf{99.1}$ \\
\midrule
\multicolumn{6}{l}{\emph{\textbf{Edge Q4\_K\_M quantization (ours)}}} \\
Q4\_K\_M PTQ & $0.851 \pm 0.011$ & $0.864$ & $0.838$ & $2.6\%$ & $91.4$ \\
Q4\_K\_M QAD & $0.911 \pm 0.008$ & $0.913$ & $0.909$ & $2.0\%$ & $97.9$ \\
\rowcolor[gray]{0.93} \textbf{Q4\_K\_M QAD + OV-Freeze} & $\mathbf{0.917 \pm 0.007}$ & $\mathbf{0.920}$ & $\mathbf{0.914}$ & $\mathbf{1.9\%}$ & $\mathbf{98.5}$ \\
\midrule
\multicolumn{6}{l}{\emph{\textbf{Domain-specific baselines}}} \\
BERT-Fraud & $0.876$ & $0.882$ & $0.870$ & $2.3\%$ & --- \\
SAFE-QAQ & $0.918 \pm 0.006$ & $0.921$ & $0.916$ & $1.8\%$ & --- \\
\bottomrule
\end{tabularx}
\end{table}

Table~\ref{tab3-en} yields several findings. First, NVFP4 QAD + OV-Freeze attains $F_1 = 0.923 \pm 0.006$, recovering 99.1\% of the BF16 upper bound ($F_1 = 0.931$) and outperforming representative PTQ methods (SpinQuant, QuaRot: $F_1 = 0.838$) by $0.085$ absolute points, as well as the self-distillation baseline BitDistiller ($F_1 = 0.858$) by $0.065$. We note that even with post-hoc temperature scaling calibration (Softmax temperature tuned to minimise KL vs.\ BF16 on 100 held-out samples), PTQ only reaches $F_1 = 0.852$ (recovery 91.5\%), remaining 7.1 points below QAD+OVF---confirming that the primary accuracy recovery stems from distillation-driven distribution alignment rather than simple temperature rescaling. Second, the QAD approach ($F_1 = 0.916$, recovery 98.4\%) substantially outperforms QAT ($F_1 = 0.844$, recovery 90.7\%) by 7.2 points, consistent with~\cite{3}'s observation that hard-label supervision degrades aligned models' refusal boundaries. Third, OV-Freeze contributes an additional $+0.7$ recovery points ($98.4\%\!\rightarrow\!99.1\%$) and reduces FPR to 1.8\%. Finally, the edge track (Q4\_K\_M QAD + OV-Freeze, $F_1 = 0.917$, recovery 98.5\%) trails the cloud scheme by only 0.006 $F_1$, confirming cross-format portability; against SAFE-QAQ ($F_1 = 0.918$, $\sim$7B params, BF16 precision, unquantized)~\cite{2}, our 248\,MB model achieves comparable accuracy with $\sim$28$\times$ storage reduction.

\medskip
\noindent\textbf{CoT contribution ablation.} Table~\ref{tab:cot-ablation-en} isolates the contribution of chain-of-thought reasoning by comparing the default pipeline against a variant that performs direct risk-score fusion without CoT. Removing CoT reduces $F_1$ by 1.8 percentage points ($0.923\!\rightarrow\!0.905$) and increases FPR from 1.8\% to 2.6\%, indicating that domain-constrained CoT provides a measurable benefit at the 0.5B scale. The degradation is most pronounced on adversarial samples (AdvFraud-3k: $F_1$ drop of 2.3 points), where multi-step semantic analysis proves most valuable.

\begin{table}[htbp]
\centering
\caption{Impact of chain-of-thought (CoT) reasoning in the cloud Tier-2 pipeline (NVFP4 QAD + OV-Freeze backbone).}
\label{tab:cot-ablation-en}
\small
\renewcommand{\arraystretch}{1.2}
\begin{tabular}{lccc}
\toprule
\textbf{Configuration} & \textbf{TAF-28k $F_1$} & \textbf{AdvFraud-3k $F_1$} & \textbf{FPR} \\
\midrule
With CoT (default) & $\mathbf{0.923 \pm 0.006}$ & $\mathbf{0.875}$ & $\mathbf{1.8\%}$ \\
Without CoT (direct fusion only) & $0.905 \pm 0.008$ & $0.852$ & $2.6\%$ \\
\bottomrule
\end{tabular}
\end{table}

\subsection{Cross-dataset Generalisation and Robustness}

\begin{table}[htbp]
\centering
\caption{Cross-dataset and adversarial robustness results}
\label{tab4-en}
\small
\renewcommand{\arraystretch}{1.15}
\begin{tabular}{>{\raggedright\arraybackslash}p{0.12\linewidth}>{\raggedright\arraybackslash}p{0.32\linewidth}>{\centering\arraybackslash}p{0.08\linewidth}>{\centering\arraybackslash}p{0.1\linewidth}>{\centering\arraybackslash}p{0.1\linewidth}>{\centering\arraybackslash}p{0.1\linewidth}}
\toprule
\textbf{Train} & \textbf{Test (property)} & \textbf{BF16} & \textbf{NVFP4 PTQ} & \textbf{NVFP4 QAD+OVF} & \textbf{$\Delta F_1$ vs BF16} \\
\midrule
TAF-28k & TAF-28k (IID) & $0.931$ & $0.838$ & $\mathbf{0.923}$ & $-0.008$ \\
TAF-28k & AdvFraud-3k (adversarial) & $0.882$ & $0.778$ & $\mathbf{0.875}$ & $-0.007$ \\
TAF-28k & ChiFraud (OOD) & $0.871$ & $0.768$ & $\mathbf{0.860}$ & $-0.011$ \\
TAF-28k + ChiFraud & AdvFraud-3k (multi-source adv.) & $0.889$ & $0.789$ & $\mathbf{0.881}$ & $-0.008$ \\
\bottomrule
\end{tabular}
\end{table}

Table~\ref{tab4-en} shows that NVFP4 QAD + OVF maintains robust classification under non-ideal distributions. On the OOD ChiFraud test set, our scheme attains $F_1 = 0.860$, a 1.1\% scalar drop vs.\ BF16 upper bound, validating cross-domain generalisation.

On the human-adversarial AdvFraud-3k, our scheme reaches $F_1 = 0.875$. Drop quantification: absolute drop vs.\ BF16 AdvFraud-3k baseline ($F_1 = 0.882$) is $0.007$; relative drop vs.\ the TAF-28k BF16 upper bound ($F_1 = 0.931$) is $(0.931 - 0.875)/0.931 = 6.0\%$ (absolute $0.056$); relative to our deployed model's TAF-28k $F_1 = 0.923$, the drop is $(0.923 - 0.875)/0.923 = 5.2\%$ (absolute $0.048$). Both relative-drop figures fall within the threat-model constraint $C_2(\le 6\%)$. Joint multi-source training (TAF-28k + ChiFraud) further lifts adversarial $F_1$ to $0.881$.

\subsection{Loss Function Ablation}

\begin{table}[htbp]
\centering
\caption{Ablation on QAD alignment loss}
\label{tab5-en}
\small
\renewcommand{\arraystretch}{1.25}
\begin{tabular}{lccc}
\toprule
\textbf{Loss variant} & \textbf{$F_1$} & \textbf{KL div.\ vs BF16} & \textbf{Recovery (\%)} \\
\midrule
Pure KL (ours) & $\mathbf{0.916 \pm 0.007}$ & $\mathbf{0.005}$ & $\mathbf{98.4}$ \\
Logits MSE & $0.901 \pm 0.010$ & $0.082$ & $96.8$ \\
Cross-entropy (NVFP4 QAT) & $0.844 \pm 0.014$ & $0.311$ & $90.7$ \\
Three-term mixed loss & $0.879 \pm 0.012$ & $0.124$ & $94.4$ \\
KL + task regularization & $0.908 \pm 0.009$ & $0.041$ & $97.5$ \\
\bottomrule
\end{tabular}
\end{table}

Pure KL achieves the Pareto optimum: highest $F_1$ and recovery with teacher-student KL compressed to 0.005. Three-term mixed loss improves over QAT only by $+3.7\%$ recovery, remaining $4.0\%$ behind pure KL, confirming that soft-label probability alignment dominates hard-label constraints ($\text{KL}\succ \text{MSE}$).

\medskip
\noindent\textbf{Interpretation of ``KL + task regularization.''} This variant appends a cross-entropy term on the task labels (fraud vs.\ benign) to the KL objective, following a two-stage schedule: pure KL for the first 1{,}400 steps, then joint KL+CE for the final 600 steps. The counterintuitive degradation (0.916$\rightarrow$0.908) arises because the teacher's soft distribution is already well-calibrated for the binary fraud-detection task (teacher $F_1 = 0.931$). Introducing hard CE labels in the late distillation phase pulls the student's non-top-1 probabilities toward a degenerate one-hot target, increasing KL divergence against the teacher (0.005$\rightarrow$0.041) and partially undoing the distributional fidelity built during pure-KL pre-training. This finding reinforces a key insight from~\cite{3}: once the teacher output distribution closely approximates the task posterior, any hard-label auxiliary loss introduces a gradient conflict with the KL objective on non-maximum tokens, particularly under low-bit quantization constraints where the student's limited capacity cannot simultaneously satisfy both signals.

\subsection{Teacher Selection Ablation}

\begin{table}[htbp]
\centering
\caption{Teacher knowledge-transfer ablation}
\label{tab6-en}
\small
\renewcommand{\arraystretch}{1.15}
\begin{tabular}{l>{\centering\arraybackslash}p{0.15\linewidth}>{\centering\arraybackslash}p{0.17\linewidth}>{\centering\arraybackslash}p{0.17\linewidth}>{\centering\arraybackslash}p{0.15\linewidth}}
\toprule
\textbf{Teacher} & \textbf{$F_1$ (fixed 0.5B tokens)} & \textbf{$F_1$ (converged)} & \textbf{Tokens to converge} & \textbf{Recovery (\%)} \\
\midrule
0.5B BF16 (homologous, ours) & $\mathbf{0.916}$ & $\mathbf{0.916}$ & $\mathbf{0.5\text{B}}$ & $\mathbf{98.4} / \mathbf{98.4}$ \\
1.8B BF16 & $0.911$ & $0.913$ & $0.7\text{B}$ & $97.9 / 98.1$ \\
3B BF16 & $0.904$ & $0.910$ & $1.0\text{B}$ & $97.1 / 97.7$ \\
7B BF16 & $0.892$ & $0.915$ & $2.0\text{B}$ & $95.8 / 98.3$ \\
\bottomrule
\end{tabular}
\end{table}

The homologous 0.5B teacher achieves the highest $F_1$ with the least data (0.5B tokens). A 7B teacher requires $4\times$ the data (2.0B tokens) for comparable fidelity, validating that low-bit QAD strongly favours homologous topologies~\cite{3}.

\subsection{OV-Freeze Ablation}

\begin{table}[htbp]
\centering
\caption{OV-Freeze layer-placement ablation}
\label{tab:ovf_layer_ablation-en}
\small
\renewcommand{\arraystretch}{1.15}
\begin{tabular}{>{\raggedright\arraybackslash}p{0.42\linewidth}cccc}
\toprule
\textbf{Strategy} & \textbf{$F_1$} & \textbf{PPL} & \textbf{$\Delta$PPL} & \textbf{Var. drift} \\
\midrule
NVFP4 QAD (no OVF baseline) & $0.916$ & $8.73$ & --- & $+18.2\%$ \\
+ OVF (FFN only) & $0.918$ & $8.71$ & $-0.02$ & $+15.4\%$ \\
+ OVF (attn $q$-proj only) & $0.918$ & $8.69$ & $-0.04$ & $+9.4\%$ \\
+ OVF (attn $q, v$) & $0.920$ & $8.66$ & $-0.07$ & $+5.1\%$ \\
+ OVF (attn $q, k, v$) & $0.922$ & $8.64$ & $-0.09$ & $+2.8\%$ \\
+ OVF (attn $q, k, v, o$) & $\mathbf{0.923}$ & $\mathbf{8.62}$ & $\mathbf{-0.11}$ & $\mathbf{+1.3\%}$ \\
+ OVF ($q,k,v,o$ + FFN) & $0.922$ & $8.63$ & $-0.10$ & $+1.5\%$ \\
\bottomrule
\end{tabular}
\end{table}

Full attention-projection OV-Freeze ($q, k, v, o$) minimises PPL ($8.62$) and variance drift ($+18.2\%\!\rightarrow\!+1.3\%$); FFN addition yields no further gain.

\begin{table}[htbp]
\centering
\caption{OV-Freeze variance estimator ablation}
\label{tab:ovf_var_ablation-en}
\small
\renewcommand{\arraystretch}{1.15}
\begin{tabular}{lcc>{\centering\arraybackslash}p{0.23\linewidth}}
\toprule
\textbf{$\mathrm{Var}(y_\ell)$ estimator} & \textbf{$F_1$} & \textbf{Pre-500-step PPL amp.} & \textbf{Convergence stability} \\
\midrule
Batch-level (instantaneous) & $0.918 \pm 0.011$ & $+1.4$ & Frequent gradient oscillation \\
EMA, $\rho = 0.95$ (default) & $\mathbf{0.923 \pm 0.006}$ & $\mathbf{+0.18}$ & Deterministic monotone \\
EMA, $\rho = 0.90$ & $0.922 \pm 0.007$ & $+0.31$ & Stable convergence \\
EMA, $\rho = 0.99$ & $0.921 \pm 0.008$ & $+0.12$ & Notably slower convergence \\
Global static prior & $0.919 \pm 0.008$ & $+0.05$ & Teacher-student drift bias \\
\bottomrule
\end{tabular}
\end{table}

EMA ($\rho = 0.95$) suppresses PPL fluctuation to $+0.18$ with deterministic monotone convergence; global static estimation lags by $0.4\%$ in final $F_1$.

\begin{table}[htbp]
\centering
\caption{OV-Freeze activation step-ratio ablation}
\label{tab9-en}
\small
\renewcommand{\arraystretch}{1.15}
\begin{tabular}{>{\raggedright\arraybackslash}p{0.3\linewidth}>{\centering\arraybackslash}p{0.15\linewidth}>{\centering\arraybackslash}p{0.15\linewidth}>{\centering\arraybackslash}p{0.29\linewidth}}
\toprule
\textbf{Activation ratio} & \textbf{$F_1$} & \textbf{PPL} & \textbf{Optimiser stability} \\
\midrule
$0\%$ (no OVF baseline) & $0.916$ & $8.73$ & Smooth \\
$10\%$ & $0.919$ & $8.68$ & Smooth \\
$20\%$ & $0.921$ & $8.65$ & Smooth \\
$\mathbf{30\%}$ \textbf{(default)} & $\mathbf{0.923}$ & $\mathbf{8.62}$ & \textbf{Smooth} \\
$40\%$ & $0.922$ & $8.63$ & Smooth \\
$50\%$ & $0.918$ & $8.66$ & Severe gradient oscillation \\
\bottomrule
\end{tabular}
\end{table}

The final 30\% activation window is Pareto-optimal: $\le 20\%$ suffers from inadequate correction; $\ge 50\%$ creates gradient conflict with the KL main loss. This matches Theorem~A.2's Lipschitz-convergence analysis.

\subsection{On-device Speculative Decoding}

\begin{table}[htbp]
\centering
\caption{Speculative decoding speedup across platforms}
\label{tab:speculative_decoding-en}
\small
\renewcommand{\arraystretch}{1.2}
\begin{tabular}{l>{\centering\arraybackslash}p{0.1\linewidth}>{\centering\arraybackslash}p{0.13\linewidth}>{\centering\arraybackslash}p{0.11\linewidth}>{\centering\arraybackslash}p{0.1\linewidth}>{\centering\arraybackslash}p{0.1\linewidth}}
\toprule
\textbf{Draft model} & \textbf{$\gamma$} & \textbf{Accept rate $\alpha$} & \textbf{Theoretical speedup} & \textbf{H100 speedup} & \textbf{SD8G3 speedup} \\
\midrule
Generic Qwen2-0.1B & $5$ & $0.78$ & $3.52\times$ & $2.92\times$ & $2.78\times$ \\
Domain-tuned (ours) & $\mathbf{5}$ & $\mathbf{0.86}$ & $\mathbf{4.25\times}$ & $\mathbf{3.49\times}$ & $\mathbf{3.32\times}$ \\
Domain-tuned & $3$ & $0.86$ & $3.24\times$ & $2.65\times$ & $2.52\times$ \\
Domain-tuned & $7$ & $0.86$ & $5.01\times$ & $4.10\times$ & $3.90\times$ \\
Domain-tuned & $10$ & $0.86$ & $5.78\times$ & $4.74\times$ & $4.51\times$ \\
\bottomrule
\end{tabular}
\end{table}

On SD8G3 at $\gamma = 5$, the system achieves $3.32\times$ end-to-end speedup. Expanding $\gamma$ to 7 yields $\sim$17\% throughput gain but 40\% KV-cache memory growth; hence $\gamma = 5$ is the static Pareto optimum on memory-constrained mobile hardware.

\medskip
\noindent\textbf{Theoretical versus measured speedup.} The theoretical $4.25\times$ (Eq.~13, $\alpha=0.86$, $\gamma=5$) exceeds measured $3.49\times$ (H100) and $3.32\times$ (SD8G3) by 18\%--22\%, due to three sources included in our measurements: (i) draft-model inference ($\sim$1.2\,ms on H100 to $\sim$4.8\,ms on SD8G3 per forward pass); (ii) KV-cache management across draft and verification passes; (iii) parallel verification dispatch ($\sim$2--3\,ms per iteration on SD8G3's heterogeneous NPU+CPU units). The measured speedups follow $\text{Speedup}_{\text{measured}} \approx \eta_{\text{plat}}\cdot (1-\alpha^{\gamma+1})/(1-\alpha)$, with $\eta_{\text{H100}} \approx 0.82$ and $\eta_{\text{SD8G3}} \approx 0.78$, stable across $\gamma$ (max deviation $\pm 0.02$).

\subsection{GLO Privacy Verification}

\begin{table}[htbp]
\centering
\caption{Speech reconstruction quality under dual white-box GLO and black-box model-inversion attacks}
\label{tab:privacy_attack-en}
\small
\renewcommand{\arraystretch}{1.25}
\begin{tabular}{l>{\centering\arraybackslash}p{0.17\linewidth}>{\centering\arraybackslash}p{0.16\linewidth}>{\centering\arraybackslash}p{0.16\linewidth}}
\toprule
\textbf{Reconstruction quality} & \textbf{White-box GLO} & \textbf{Black-box U-Net inversion} & \textbf{Random baseline} \\
\midrule
PESQ (1--5) & $1.21$ & $1.16$ & $1.05$ \\
WER & $\ge 0.95$ & $\ge 0.97$ & $1.00$ \\
MOS (1--5) & $1.18$ & $1.11$ & $1.00$ \\
Speaker-ID accuracy & $8.3\%$ & $7.9\%$ & $10.0\%$ \\
$\hat{I}(x; F_v)$ & $\approx 0$ & $\approx 0$ & --- \\
\bottomrule
\end{tabular}
\end{table}

Under both high-strength attacks, reconstructed-speech WER approaches the complete-collapse upper bound ($\ge 0.95$), indicating that the embedding compression renders the waveform practically unreconstructable. Speaker-ID accuracy of 7.9--8.3\% is below the 10\% statistical random baseline. This confirms the empirical non-invertibility of $F_v$ at both semantic and speaker-identity levels, satisfying PIPL compliance and the threat-model constraint $\mathcal{G}_1$.

% ============================================================================
% 5. DISCUSSION
% ============================================================================
\section{Discussion and Limitations}

\subsection{Multimodal-detection mechanism}

QAD-MultiGuard's $F_1 = 0.923$ on TAF-28k stems from low-bit quantization compactness and cross-modal complementarity: acoustic paralinguistic features (speech-rate anomalies, anxiety tones) compensate where text-only signals are ambiguous (e.g., colloquial impersonation scripts). OV-Freeze maintains close alignment with the BF16 teacher (KL $\le 0.005$).

\subsection{Privacy compliance}

Raw audio is feature-extracted on-device and then discarded; only temporally-compressed non-invertible $F_v$ is transmitted. White-box GLO attack failure (PESQ 1.21, WER $\ge 0.95$) confirms reverse-recovery infeasibility. This feature-discard mechanism avoids DP noise injection's accuracy penalties while providing a quantifiable PIPL compliance baseline.

\subsection{Misclassification analysis and limitations}

Error tracing of 213 misclassified samples reveals three limitations: (i) $\sim$41\% involve audio $<8$\,s, where temporal context is insufficient; (ii) white-box PGD-20 attack ($\epsilon = 0.10$) causes monotone accuracy degradation under deliberate perturbations; (iii) extreme noise (SNR $<5$\,dB) induces $\sim$4.2\% discrimination drop.

\subsection{Distribution-shift assessment}

TAF-28k is essentially a ``re-enactment'' lab distribution. Generalisation testing on 54 real law-enforcement case recordings showed an $F_1$ drop of $\sim$6.8\%, attributable to channel loss, regional dialects, and adversarial script adaptation in real fraud calls. This gap suggests the reported figures represent the architecture's empirical upper bound under the current setup, rather than its expected performance in operational law-enforcement deployment. Once Blackwell hardware is broadly available for production deployment, we plan to replace the current NBE emulation data with native NVFP4 measurements on actual hardware.

% ============================================================================
% 6. CONCLUSION
% ============================================================================
\section{Conclusion}

This paper presents QAD-MultiGuard, an edge-side multimodal anti-fraud framework integrating pure-KL quantization-aware distillation, OV-Freeze regularization, and a 128-d non-invertible acoustic embedding---all within a 240\,MB on-device footprint. On TAF-28k, it achieves $F_1 = 0.923$, recovering 99.1\% of the BF16 upper bound; with domain-adapted speculative decoding ($\gamma = 5$, $\alpha = 0.86$), SD8G3 attains $3.32\times$ speedup. Dual GLO and model-inversion attacks confirm $F_v$ non-invertibility (WER $\ge 0.95$).

Limitations include short-call robustness ($<$8\,s audio) and real-world distribution shift ($\sim$6.8\% $F_1$ drop on forensic recordings). Future work will pursue multimodal adversarial training, long-temporal modelling for short calls, and semi-supervised domain adaptation to bridge the lab-to-field gap.

% ============================================================================
% APPENDIX A — FORMAL ANALYSIS
% ============================================================================
\appendix
\section*{Appendix A: Formal Analysis of OV-Freeze and Acoustic Embedding Non-invertibility}
\addcontentsline{toc}{section}{Appendix A: Formal Analysis of OV-Freeze and Acoustic Embedding Non-invertibility}

\subsection*{A.1 Proposition 1: Boundedness of the OV-Freeze correction factor}

\noindent\textbf{Proposition A.1 (Boundedness of $c_\ell$).} Let the original output of the quantization-sensitive layer $\ell$ be $\bm{y}_\ell\in\mathbb{R}^{N\times d}$, and the OV-Freeze-corrected output be $\bm{y}'_\ell = c_\ell\,\bm{y}_\ell$, where the scalar correction factor is
\[
c_\ell = \sqrt{\frac{\sigma^2_{\mathrm{BF16},\ell}}{\mathrm{Var}(\bm{y}_\ell)+\epsilon}},\qquad \epsilon = 10^{-6}.
\]
Then $c_\ell$ lies in a compact positive interval $[c_{\min}, c_{\max}] \subset (0, \infty)$ for any finite batch output $\bm{y}_\ell$, bounded away from zero and infinity.

\noindent\textbf{Proof.} The teacher baseline $\sigma^2_{\mathrm{BF16},\ell}>0$ is a fixed positive constant. The empirical variance $\mathrm{Var}(\bm{y}_\ell)$ is finite for any finite batch ($\mathrm{Var}(\bm{y}_\ell) < \infty$). Since $\epsilon>0$, the denominator $\sqrt{\mathrm{Var}(\bm{y}_\ell)+\epsilon}$ is finite and strictly positive; thus $c_\ell$ is bounded above by $\sigma_{\mathrm{BF16},\ell}/\sqrt{\epsilon}$ and below by $\sigma_{\mathrm{BF16},\ell}/\sqrt{\mathrm{Var}_{\max}+\epsilon}$, where $\mathrm{Var}_{\max}$ is the maximal variance observed across training (bounded by the BF16 teacher's variance scale). Therefore $c_\ell \in [c_{\min}, c_{\max}]$.\hfill$\square$

\medskip
\noindent\textbf{Connection to gradient compatibility.} In the operational forward--backward graph (\S~3.2.5, Eq.~\eqref{eq:ovf-detach-en}), the variance branch is detached, making $c_\ell$ a constant scaling factor. The gradient is thus $\partial L_{\mathrm{QAD}} / \partial \bm{y}_\ell = c_\ell \cdot \partial L_{\mathrm{QAD}} / \partial \bm{y}'_\ell$ (Eq.~\eqref{eq:ovf-detach-en}), which inherits the Lipschitz properties of the KL loss scaled by a bounded constant $c_\ell \in [c_{\min}, c_{\max}]$. For completeness, the full undetached gradient (Eq.~\eqref{eq:ovf-full-grad-en}) would include an additional Jacobian term $\partial c_\ell / \partial \bm{y}_\ell$:
\[
\frac{\partial L_{\mathrm{QAD}}}{\partial\bm{y}_\ell} = c_\ell\cdot\frac{\partial L_{\mathrm{QAD}}}{\partial\bm{y}'_\ell} + \left\langle\frac{\partial L_{\mathrm{QAD}}}{\partial\bm{y}'_\ell},\,\bm{y}_\ell\right\rangle\cdot\frac{\partial c_\ell}{\partial\bm{y}_\ell},
\]
which under the full backprop path (without \texttt{detach}) satisfies the strict bound
\[
\left\|\frac{\partial c_\ell}{\partial \bm{y}_\ell}\right\|_2 \;\le\; \frac{\sigma_{\mathrm{BF16},\ell}}{2\bigl(\mathrm{Var}(\bm{y}_\ell)+\epsilon\bigr)^{3/2}}\cdot\frac{2\|\bm{y}_\ell\|_2}{N-1},
\]
ruling out gradient divergence even if the detachment were removed. The practical detached version is therefore a conservative simplification that strictly reduces gradient complexity while inheriting all convergence guarantees.

\subsection*{A.2 Theorem 1: Lipschitz constant and SGD convergence of the joint loss}

\noindent\textbf{Setup.} Let $L_{\mathrm{KL}}$ denote the KL main loss, with Lipschitz constant $L$ under BF16 teacher variance calibration. OV-Freeze multiplies each sensitive projection output by $c_\ell$ in the final $30\%$ of training steps. Since the EMA-estimated $\mathrm{Var}^{(t)}_{\mathrm{EMA}}(\bm{y}_\ell)$ is initialised to $\sigma^2_{\mathrm{BF16},\ell}$, $c_\ell \in [c_{\min}, c_{\max}]$ falls in a compact interval; denote $c_{\max} = \max_\ell \|c_\ell\|_\infty$.

\noindent\textbf{Theorem A.2 (Joint-loss convergence).} If $L_{\mathrm{KL}}$ has Lipschitz-continuous gradient with constant $L$ on parameter space $\Theta$, and the learning rate $\eta \le 1/L'$ where
\begin{equation}
\label{eq:Lprime-en}
L' = L\cdot\max_\ell\|c_\ell\|_\infty,
\end{equation}
then the SGD iterate $\{\theta_t\}$ satisfies
\begin{equation}
\label{eq:conv-en}
\min_{t \le T}\;\mathbb{E}\!\left[\|\nabla L_{\mathrm{joint}}(\theta_t)\|^2\right] \;\le\; \frac{2(L_{\mathrm{joint}}(\theta_0) - L^\star)}{\eta\,T}.
\end{equation}

\noindent\textbf{Proof sketch.} Under the detached gradient (Eq.~\eqref{eq:ovf-detach-en}), OV-Freeze multiplies the KL gradient by a per-layer constant $c_\ell$ bounded in $[c_{\min}, c_{\max}]$ (Proposition~A.1). The effective loss is $L_{\mathrm{joint}} = c_\ell \cdot L_{\mathrm{KL}}$, whose Lipschitz constant is $L' = L \cdot \max_\ell\|c_\ell\|_\infty$ by the chain rule for scaled functions. Substituting into the standard SGD convergence analysis (Bubeck, \textit{Convex Optimization: Algorithms and Complexity}, 2015, Theorem~3.7) yields \eqref{eq:conv-en}.\hfill$\square$

\noindent\textbf{Discussion.} Eq.~\eqref{eq:Lprime-en} shows that $L'$ grows monotonically with $\max_\ell\|c_\ell\|_\infty$. Even under detachment, $c_\ell$ is not fixed across training---it depends on the ratio $\sigma_{\mathrm{BF16},\ell} / \sqrt{\mathrm{Var}_{\mathrm{EMA}}(\bm{y}_\ell)}$, and the EMA estimate evolves as training progresses. In early training (first $\sim$70\% of steps, during which OV-Freeze is not yet active), $\mathrm{Var}(\bm{y}_\ell)$ may still be far from the teacher baseline. When OV-Freeze activates at the final 30\% window, if the activation ratio were $\ge 50\%$, the learning rate $\eta = 1\times 10^{-5}$ may no longer satisfy $\eta \le 1/L'$, causing gradient oscillation. The 30\% window empirically ensures that when OV-Freeze activates, $\mathrm{Var}_{\mathrm{EMA}}(\bm{y}_\ell) \approx \sigma^2_{\mathrm{BF16},\ell}$ and thus $c_\ell \approx 1$, keeping $L' \approx L$ and $\eta L' \ll 1$. This explains why the detached implementation still exhibits sensitivity to the activation ratio (Table~\ref{tab9-en}): $c_\ell$ converges to 1 only within the late-training window, and a premature activation start causes $c_\ell$ to act as a large constant multiplier of the KL gradient, destabilising the optimisation trajectory.

\subsection*{A.3 Proposition 2: Engineering non-invertibility guarantee for $F_v$}

\noindent\textbf{Proposition A.3 (Non-invertibility guarantee via dual temporal compression).} Let the raw PCM signal be $x\in\mathbb{R}^{S}$, and the acoustic embedding $F_v\in\mathbb{R}^{128}$ be the concatenation of 64-d time-averaged MFCC and 64-d Whisper-tiny encoder global-pooled projection features. Then $F_v$ is empirically non-invertible: under white-box GLO and black-box model-inversion attacks, reconstructed speech quality satisfies $\text{WER}\ge 0.95$, $\text{PESQ}\le 1.21$, and speaker-ID accuracy $\le 8.3\%$ (approaching the 10\% random baseline).

\noindent\textbf{Rationale.} Non-invertibility arises from two independent temporal compression operations, neither of which preserves frame-level phase or temporal structure:

\begin{itemize}
\item \textbf{MFCC branch:} Frame-level MFCC features (64-d, extracted at 10\,ms hop, producing $\sim$300 frames per 3\,s window) are mean-pooled along the time dimension into a single 64-d vector $\bm{f}_{\mathrm{mfcc}}$. This averaging collapses $\mathcal{O}(300)$ frames of phoneme-level dynamics into static statistical moments, discarding temporal ordering, phase, and prosodic timing information essential for waveform reconstruction.

\item \textbf{Whisper branch:} The Whisper-tiny encoder outputs $T$ frame-level features $\in\mathbb{R}^{T\times 384}$, which are globally average-pooled to $\bar{\bm{h}}_w\in\mathbb{R}^{384}$ and then projected via $\bm{W}_{\mathrm{proj}}\in\mathbb{R}^{64\times 384}$ to 64-d. Since $\mathrm{rank}(\bm{W}_{\mathrm{proj}})\le 64$, this rank-limited projection further compresses the pooled representation. The sequential combination of global pooling (discarding temporal ordering) and low-rank projection (discarding nullspace information) severs the mapping from $F_v$ back to the original waveform at the physical compression level.
\end{itemize}

\noindent\textbf{In contrast to information-theoretic non-invertibility.} We note that this is an \emph{engineering} non-invertibility guarantee validated by adversarial reconstruction experiments, rather than a formal proof that $I(x;F_v)=0$. Because $F_v$ retains task-relevant statistical properties of the acoustic signal (e.g., energy variance, pitch range), a small but non-zero mutual information $I(x;F_v)>0$ necessarily exists. The empirical evidence---WER~$\ge 0.95$, PESW~$\le 1.21$, speaker-ID accuracy approaching random baseline---demonstrates that this residual information is insufficient for any known inverse attack to recover intelligible speech or identifiable speaker characteristics, satisfying PIPL compliance requirements for practical deployment.

\medskip
\noindent\textbf{Empirical validation.} Table~\ref{tab:privacy_attack-en} in \S~4.10 confirms the non-invertibility at both semantic and speaker-identity levels: under both white-box GLO and black-box model-inversion attacks, reconstructed speech yields WER~$\ge 0.95$, PESQ~$\le 1.21$, and speaker-ID accuracy~$\le 8.3\%$. The $\epsilon$-LDP extension (\S~3.3.2, Eq.~12) provides an additional formal privacy guarantee for extreme-privacy scenarios beyond the baseline engineering guarantee.\hfill$\square$

\newpage
\printbibliography
\end{document}
